% Copy each caption into the corresponding figure environment.
% Keep your existing labels and cross-references.

% Figure 1. Performance of forecasting, nowcasting and contemporaneous models
\caption{\textbf{Performance of forecasting, nowcasting and contemporaneous models.} \textbf{a,} Grey bars show exact overall-phase classification accuracy. \textbf{b,} Sensitivity (recall) and precision for identifying Integrated Food Security Phase Classification (IPC) Phase 3 or higher: blue circles denote forecasting, orange squares nowcasting and green triangles contemporaneous modelling. \textbf{c,} Predicted versus observed population shares in IPC Phase 3 or higher. Orange circles represent observations; overlapping semi-transparent circles appear darker. Solid vermilion lines are least-squares regressions of observed on predicted shares, and dashed grey lines indicate equality (y = x, where x is the prediction and y the observation). R$^2$ is the coefficient of determination calculated from the original predictions, not the fitted line. Forecasting and nowcasting use the 2022 temporal holdout (n = 1,170 observations each), while contemporaneous results use pooled random five-fold, row-level out-of-fold predictions (n = 5,575; seed 0). Evaluation populations and validation protocols differ across models. Source data are provided as a Source Data file.}

% Figure 2. SHAP contributions of forecasting and nowcasting features
\caption{\textbf{SHAP contributions of forecasting and nowcasting features.} \textbf{a, b,} Ten leading individual features for forecasting and nowcasting layer 2, respectively, with Phase 3 or higher above Phase 4 or higher. \textbf{c, d,} Corresponding contributions of the selected variable groups. SHAP denotes SHapley Additive exPlanations. Individual importance is the absolute value of the signed SHAP sum across 4,405 training observations from 2017–2021. Group importance sums signed contributions within each original variable group and across observations, takes the absolute value, and normalizes across the included groups. These summaries are not mean absolute SHAP values. Nowcasting layer 2 explains the residual correction to layer 1. Green denotes Phase 3 or higher and blue Phase 4 or higher. Phase 3+ and Phase 4+ refer to the corresponding cumulative Integrated Food Security Phase Classification (IPC) categories. In feature labels, a suffix + denotes a 12-month lag; * denotes the 12-month summary used in the original notebook (moving mean, or temporal standard deviation where labelled SD). SD means standard deviation; WFP, World Food Programme; VIIRS, Visible Infrared Imaging Radiometer Suite; GDP, gross domestic product; EVI, enhanced vegetation index; and pH, the measure of soil acidity. The five- and ten-station labels refer to the corresponding spatially aggregated conflict features. Source data are provided as a Source Data file.}

% Figure 3. Ensemble procedure for assigning an IPC phase
\caption{\textbf{Ensemble procedure for assigning an IPC phase.} Four extreme gradient boosting (XGBoost) regressors estimate cumulative population shares in Integrated Food Security Phase Classification (IPC) Phases 2+, 3+, 4+ and 5; + means the stated phase or higher. Pale-blue rectangles denote regressors, pale-green ellipses predicted population shares, grey diamonds threshold decisions and blue output rectangles assigned phases. Arrows show information flow and sequential decisions. Starting at Phase 5, an area is assigned the highest phase whose predicted cumulative population share is at least 20\% (the symbol $\geq$ means greater than or equal to). A Yes branch assigns that phase; a No branch moves to the next lower threshold. If none qualifies, Phase 1 is assigned. Source data are provided as a Source Data file.}

% Figure 4. Distribution of recorded overall IPC phases
\caption{\textbf{Distribution of recorded overall IPC phases.} \textbf{a,} Counts across all 5,575 observations from 2017–2022. \textbf{b,} Percentage distribution within each year. IPC AFI denotes Integrated Food Security Phase Classification Acute Food Insecurity. Green indicates Phase 1 (Minimal), light blue Phase 2 (Stressed), blue Phase 3 (Crisis), orange Phase 4 (Emergency) and vermilion Phase 5 (Famine). Each stacked segment in b represents that phase's share of the year's observations. Both panels use overall-phase labels recorded in the source data. Source data are provided as a Source Data file.}

% Supplementary Figure 1. Observation frequency across analysis areas
\caption{\textbf{Observation frequency across analysis areas.} Blue bars show the number of analysis areas with each observation frequency. The sample contains 5,575 observations from 1,198 unique analysis areas. Of these, 510 areas (42.6\%) appear no more than twice; the median frequency is three observations and the maximum is 20. Each area contributes once to the histogram. Source data are provided as a Source Data file.}

% Supplementary Figure 2. Predicted and observed areas above the 70th-percentile severity threshold
% IPC usage evidence checked 5 November 2026: the IPC-published HDX country
% datasets include GeoJSON polygons and state "Public Domain / No restrictions (CC0)".
% Publisher and data licence: https://data.humdata.org/dataset/ethiopia-acute-food-insecurity-country-data
% Machine-readable metadata: https://data.humdata.org/api/3/action/package_show?id=ethiopia-acute-food-insecurity-country-data
% CC0 usage statement: https://creativecommons.org/publicdomain/zero/1.0/
% API-to-HDX provenance: https://github.com/OCHA-DAP/hdx-scraper-ipc/blob/325b1505b6ca7f52ef4aebb45b02eb51cc4fc594/src/hdx/scraper/ipc/ipc.py#L292
% Scope: these establish the current HDX data release terms, not the historical
% API terms or exact polygon-file identity on 30 April 2023. The caption keeps this distinction.
\caption{\textbf{Predicted and observed areas above the 70th-percentile severity threshold.} \textbf{a,} Forecasting, 12 months ahead. \textbf{b,} Observed outcomes. Severity is the population share in Integrated Food Security Phase Classification (IPC) Phase 3 or higher. Each of the 646 analysis areas is represented by its maximum predicted share in 2022 for a and its maximum observed share in 2022 for b. Vermilion indicates areas at or above the panel-specific 70th percentile; blue indicates the remaining evaluated areas. Cutoffs are 0.35 and 0.50, respectively. Ties are retained, yielding 207 predicted and 195 observed high-severity areas, so the selected proportion can exceed 30\%. Pale grey denotes background land, thin grey lines country boundaries, white lines analysis-area boundaries and white background water. Background country boundaries are from Natural Earth (public domain). Study-area boundaries are IPC polygons obtained through the IPC application programming interface (API; accessed 30 April 2023). IPC's current GeoJSON datasets on the Humanitarian Data Exchange (HDX) are designated Public Domain/No restrictions (CC0). Source data are provided as a Source Data file.}

% Supplementary Figure 3. IPC predictions before and after combined Phase 4/5 rescue
\caption{\textbf{IPC predictions before and after combined Phase 4/5 rescue.} Rows show forecasting (a–b), nowcasting (c–d) and contemporaneous modelling (e–f). IPC denotes Integrated Food Security Phase Classification; P denotes phase; and P4/5 combines phases 4 and 5. Left panels retain five predicted classes. Right panels combine predicted phases 4 and 5 while retaining all five actual classes and the same observations. Rescue promotes eligible baseline Phase 3 predictions to combined P4/5 and does not distinguish Phase 4 from Phase 5. Each cell reports its count and percentage within the actual-phase row. The dark-blue-to-yellow colour scale encodes the same within-row share, from zero to one. Forecasting and nowcasting use the square-root-balance policy on the 2022 temporal holdout (n = 1,170 each). Contemporaneous modelling uses full balance and random five-fold, row-level out-of-fold (OOF) predictions (n = 5,575; seed 0), with the threshold selected from pooled OOF results. Actual contemporaneous phases are reconstructed using the cumulative-share 20\% rule; Figure 4 instead uses recorded source labels. Source data are provided as a Source Data file.}

% Supplementary Figure 4. Actual and predicted cumulative population shares
\caption{\textbf{Actual and predicted cumulative population shares.} Columns show forecasting, nowcasting and contemporaneous estimates. Rows show Integrated Food Security Phase Classification (IPC) Phase 2+ (a–c), Phase 3+ (d–f) and Phase 4+ (g–i); + means the stated phase or higher. Orange circles represent observations; overlapping semi-transparent circles appear darker. Solid vermilion lines regress observed shares on predicted shares, and dashed grey lines indicate equality (y = x, where x is the prediction and y the observation). An intercept of zero and slope of one are the linear-calibration reference. R$^2$ is the coefficient of determination comparing the original predictions with observations, not the fitted line's goodness of fit. These are population shares, not probabilities that an area experiences a crisis. Forecasting and nowcasting use the 2022 temporal holdout (n = 1,170 per panel); contemporaneous estimates use random five-fold, row-level out-of-fold predictions (n = 5,575 per panel). Source data are provided as a Source Data file.}

% Supplementary Figure 5. Binary confusion matrices for the Phase 4/5 rescue models
\caption{\textbf{Binary confusion matrices for the Phase 4/5 rescue models.} Rows show forecasting, nowcasting and contemporaneous modelling. Columns show unweighted (mint), square-root-balance (gold) and full-balance (salmon) policies. IPC denotes Integrated Food Security Phase Classification; P4/5 means IPC Phase 4 or 5, and below P4 means phases 1–3. Rows within each matrix are actual classes and columns predicted classes. Each cell reports its count and percentage within the actual-class row. Stronger colour represents a larger percentage, as shown by the three scales; hues identify policies. Black frames mark the primary policy selected separately for each modelling task using out-of-fold (OOF) results. Forecasting and nowcasting use the fixed 2022 temporal holdout (n = 1,170 each). Contemporaneous modelling uses random five-fold, row-level OOF predictions (n = 5,575; seed 0), with threshold selection based on pooled OOF predictions. Validation protocols and populations differ across rows. Source data are provided as a Source Data file.}

% Supplementary Figure 6. Temporal validation for forecasting hyperparameter selection
\caption{\textbf{Temporal validation for forecasting hyperparameter selection.} Blue boxes denote training years, pale-orange boxes validation years and the grey box the held-out 2022 test year. Black arrows indicate fitting on the training period followed by evaluation in the next designated year. The bracket groups the four expanding-window folds used for hyperparameter selection: training begins in 2017, with successive validation years 2018–2021. The 2022 observations are excluded from every tuning fold. For final forecasting and nowcasting evaluation, 2017–2021 form the training set and 2022 the fixed test set. Nowcasting and contemporaneous models use fixed hyperparameters without separate hyperparameter searches. Source data are provided as a Source Data file.}
